\section{Centered Orthogonal Routing for Experts}
\label{sec:method}

The analysis above leads to a simple design principle: a useful MoE gate must remove the shared global offset, but it must not discard the mean-direction information that still carries coarse routing cues. CORE implements this idea with an explicitly decoupled first-order gate.

\subsection{Routing Decomposition}
Given the hidden representation $x$, we define
$$
u=\frac{\mu_{\mathrm{global}}}{\|\mu_{\mathrm{global}}\|_2}, \qquad
x_{\parallel}=(u^\top x)u, \qquad
x_c=x-\mu_{\mathrm{global}}.
$$
The router for expert $e$ is then
$$
\ell_e(x)=a_e^\top x_{\parallel}+b_e^\top x_c,
$$
where $a_e,b_e \in \mathbb{R}^d$.

This decomposition creates two explicit routing channels:
\begin{itemize}
    \item the \textbf{mean channel} $a_e^\top x_{\parallel}$, which captures coarse variation aligned with the dominant mean-related direction;
    \item the \textbf{residual channel} $b_e^\top x_c$, which captures centered first-order variation after the shared global offset has been removed.
\end{itemize}

\subsection{Why This Form Is Necessary}
The design is motivated by two complementary failure modes.

First, uncentered routing entangles a large shared bias with content-dependent variation. This makes the gate sensitive to the dominant global direction in a way that can produce collapse-like assignments.

Second, centered routing alone is incomplete if it simply discards the mean direction. The mean-related axis still contains useful routing information, especially for function words, high-frequency tokens, and other coarse token classes. CORE therefore preserves that axis explicitly, rather than forcing a single matrix to absorb both coarse and fine routing roles.

\subsection{Specialization-Oriented Objective}
CORE is trained with the standard task loss. In the current first-order version, we do not rely on quadratic routing terms. The method claim is therefore minimal: if the representation regime is predominantly mean-separable, then expert specialization can be improved by routing with a centered and orthogonally decoupled first-order gate.

\subsection{Task-Agnostic Routing}
CORE does not use task identity at inference time. Routing is computed directly from the current hidden representation. This makes the method naturally compatible with task-agnostic continual-learning protocols, although the value of this property must still be validated in downstream experiments.

\subsection{Computational Simplicity}
The computational cost of CORE is linear in the hidden dimension per expert. Both routing channels are first-order projections, so the overhead is modest and easy to analyze. This simplicity is not just an engineering convenience; it is part of the method claim. If the current regime is already well served by decoupled first-order routing, then introducing higher-order complexity is unnecessary.
